\cotefile{22}
\scanpage{2}
% printed running head
\textbf{C. R. Acad. Sc. Paris, t. 257, p. 2948--2950 (13 novembre 1963).}\hfill\textbf{Groupe 1.}\par\noindent\rule{\linewidth}{0.4pt}

\textsc{Algèbre.} --- \textit{Étude différentielle des anneaux locaux réguliers : applications.} Note (\textsuperscript{*}) de M. \textbf{\textsc{Jean-Pierre Jouanolou}}, présentée par M. René Garnier.

Le paragraphe 1 donne pour les anneaux locaux usuels (localités, localités analytiques, etc.) une caractérisation différentielle de régularité. Remarquant que, pour ces anneaux, la complétion se fait de façon « séparable », on montre que la réduction et la normalité analytiques en résultent simplement.

1. Les critères différentiels de régularité de Nakai, Suzuki et Mount font \deleted{apparaître}\inserted{\handwritten{appel}} en caractéristique $p>0$ \inserted{\handwritten{$/$ à}} une \deleted{condition}\inserted{\handwritten{$\angle$ hypothèse}} de réduction analytique (il y a des exemples contraires sinon). L’énoncé suivant permet d’éviter cette condition dans les cas usuels :
% The two oblique proof marks at the right repeat the insertion and substitution signs.

\margin{\handwritten{$\left.\rule{0pt}{7em}\right]$}}\textsc{Proposition 1.} --- \textit{Soient $\mathrm A$ un anneau local régulier équicaractéristique tel que \tikz[baseline=(x.base)]{\node[draw,ellipse,inner sep=1pt] (x) {$\mathrm D(\mathrm A)$};} soit plat et $\mathfrak a$ un idéal de $\mathrm A$ tel $\mathrm A/\mathfrak a$ soit \tikz[baseline=(x.base)]{\node[draw,ellipse,inner sep=1pt] (x) {\itshape réduit.};}}

\textit{Les assertions suivantes sont équivalentes :}

\textit{(a) $\mathrm A/\mathfrak a$ est régulier;}

\textit{(b) $\operatorname{Tor}^{1}_{\mathrm A/\mathfrak a}(\mathrm A/\mathfrak m,\ \mathrm D(\handwritten{\not}\mathrm A))=0$;}
\margin{\handwritten{$A/\underline{a}\,/$}}
% Matching oblique reference strokes touch the A inside D(A) in (b) and (c).

\textit{(c) $\operatorname{Tor}^{i}_{\mathrm A/\mathfrak a}(\mathrm A/\mathfrak m,\ \mathrm D(\handwritten{\not}\mathrm A))=0$ pour tout $i$ supérieur ou égal à $1$ ($\mathfrak m$ désignant l’idéal maximal de $\mathrm A$).}

\textit{De plus, si $\mathrm D(\mathrm A)$ est libre, (a) équivaut à :}

\textit{(d) $\mathrm D(\mathrm A/\mathfrak a)$ est libre.}

\textsc{Corollaire 1.} --- \textit{Soient $k$ un corps valué complet de caractéristique $p>0$ et tel que $(k:k^{p})<+\infty$, $\mathrm B$ une algèbre de séries convergentes sur $k$ et $\mathrm A$ un anneau réduit localisé d’un anneau quotient de $\mathrm B$. Pour que $\mathrm A$ soit régulier, il faut et il suffit que $\mathrm D(\mathrm A)$ soit libre.}

\textsc{Corollaire 2} [Kunz (\textsuperscript{3})]. --- \textit{Pour qu’une localité \uline{réduite} $\mathrm A$ soit régulière, il faut et il suffit que son module de différentielles absolu $\mathrm D(\mathrm A)$ soit libre.}
\margin{\handwritten{$\big|$ contenu dans le critère jacobien \ldots}}

\textsc{Corollaire 3} [Suzuki (\textsuperscript{10}), Mount (\textsuperscript{5})]. --- \textit{Soit $\mathrm A$ un anneau local nœthérien analytiquement \uline{réduit} de caractéristique $p>0$ et tel que le séparé $\mathrm D'(\mathrm A)$ soit de \uline{type fini}. Alors, pour que $\mathrm D'(\mathrm A)$ soit libre, il faut et il suffit que $\mathrm A$ soit régulier.}
\margin{\handwritten{assez rare !}}

La proposition 1 fournit des critères pour que l’ensemble des points non réguliers du spectre d’un anneau (lieu singulier) soit un fermé, notamment :

\margin{\handwritten{$\left.\rule{0pt}{2em}\right]$}}\textsc{Proposition 2.} --- \textit{Si $\mathrm A$ est un anneau global régulier équicaractéristique tel que $\mathrm D(\mathrm A)$ soit libre, le lieu singulier de toute $\mathrm A$-algèbre affine $\mathrm B$ est fermé.}
\margin{\handwritten{\unsure{séries|série} \unsure{formelles|formelle} ??}}

En particulier :

\textsc{Proposition 3.} --- \textit{Soit $\mathrm G$ un schéma en groupes régulier sur un corps parfait. Alors tout schéma de type fini sur $\mathrm G$ a un lieu singulier fermé.}
\margin{\handwritten{\unsure{canularesque|canularique|canulariques}}}

2. Dans tout ce qui suit, $\mathrm A$ étant un anneau local, nous désignerons par $\mathrm L_{\mathrm A}$ son corps résiduel et par $\mathrm K_{\mathrm A}$ son corps de fractions s’il est intègre.
% archive sheet number: 1
\scanpage{3}
\mspage{( 2 )}

Soient $\mathrm A$ un anneau local nœthérien et $\widehat{\mathrm A}$ son complété. On a :

\textsc{Lemme.} --- \textit{Le noyau de l’application canonique : $\widehat{\mathrm A}\underset{\mathrm A}{\otimes}\mathrm D(\mathrm A)\to\mathrm D(\widehat{\mathrm A})$ est contenu dans l’intersection : $\displaystyle\bigcap_{\mathrm M\in\sigma}\left(\widehat{\mathrm A}\underset{\mathrm A}{\otimes}\mathrm M\right)$, où $\sigma$ désigne l’ensemble des sous-modules $\mathrm M$ de $\mathrm D(\mathrm A)$ tels que $\mathrm F=\mathrm D(\mathrm A)/\mathrm M$ soit un $\mathrm A$-module de type fini.}

De ce lemme on déduit que si $\mathrm D(\mathrm A)$ est semi-fini au sens de Suzuki (\textsuperscript{10}), c’est-à-dire somme directe d’un module libre et d’un module de type fini, la suite canonique
\[
0\to\widehat{\mathrm A}\underset{\mathrm A}{\otimes}\mathrm D(\mathrm A)\to(\widehat{\mathrm A})\to\mathrm D_{\mathrm A}(\widehat{\mathrm A})\to0
\]
est exacte.

Cela entraîne notamment :

\textsc{Proposition 4.} --- \textit{Si $\mathrm A$ est une localité analytique\uline{ment intègre}, $\mathrm K_{\widehat{\mathrm A}}$ est séparable sur $\mathrm K_{\mathrm A}$.}

De même (\textit{cf.} cor. 1 de la prop. 1), si $\mathrm A$ est un anneau de séries convergentes sur un corps valué complet de caractéristique $p>0$ et tel que $(k:k^p)<+\infty$, $\mathrm K_{\widehat{\mathrm A}}$ est séparable sur $\mathrm K_{\mathrm A}$. Ceci permet de combler ce qui nous paraît être une lacune dans la démonstration donnée par Lang (\textsuperscript{4}) du résultat suivant :
% slanted
\margin{\handwritten{$\left.\rule{0pt}{8em}\right]$ connu}}

\textsc{Proposition 5.} --- \textit{Si $\mathrm A=k\{\mathrm T\}$ est l’anneau des séries convergentes sur un corps $k$ valué complet et algébriquement clos, $\mathrm K_{\mathrm A}$ est quasialgébriquement clos $(\mathrm C^1)$.}

Un cas particulier de la proposition 4, classique, est que $\mathrm K((\mathrm X_1,\mathrm X_2,\ldots,\mathrm X_n))$ est séparable sur $\mathrm K(\mathrm X_1,\ldots,\mathrm X_n)$ pour tout corps $\mathrm K$. Il implique, en tenant compte du théorème de structure de Cohen :

\textsc{Proposition 6.} --- \textit{Soit $\mathrm A$ un anneau local régulier contenant un corps $k$ tel que $\mathrm L_{\mathrm A}$ soit séparable sur $k$. Alors $\mathrm K_{\mathrm A}$ est aussi séparable sur $k$.}

3. \textsc{Critère de réduction analytique.} --- La comparaison de la proposition 1 et des énoncés de Suzuki et Mount suggère l’existence d’un critère de réduction analytique sous-jacent; nous nous proposons de voir dans ce paragraphe comment les conclusions du paragraphe 2 permettent d’en obtenir un. Notons tout d’abord le lemme suivant qui peut être considéré comme classique :

\textsc{Lemme.} --- \textit{Soit $\mathrm A$ un anneau local nœthérien analytiquement intègre et tel que $\mathrm K_{\widehat{\mathrm A}}$ soit séparable sur $\mathrm K_{\mathrm A}$. Si $\mathrm B$ est un anneau local intègre de la forme $\mathrm B=\mathrm C_{\mathfrak p}$, où $\mathrm C$ est une $\mathrm A$-algèbre \uline{finie} et $\mathfrak p$ un idéal premier de $\mathrm C$ tel que $\mathfrak p\cap\mathrm A=\mathfrak M$ ($\mathfrak M$ idéal maximal de $\mathrm A$), $\mathrm B$ est analytiquement réduit.}
% slanted
\margin{\handwritten{$\big|$ faux si on ne \unsure{suppose} pas $A\to B$ injectif ??}}

On en déduit :

\textsc{Théorème.} --- \textit{Soit $\mathrm A$ un anneau local régulier vérifiant l’une ou l’autre des conditions suivantes :}

$(h_1)$ \textit{$\mathrm K_{\mathrm A}$ est de caractéristique $0$;}

$(h_2)$ \tikz[baseline=(x.base)]{\node[draw,ellipse,inner sep=2pt] (x) {\itshape $\mathrm D(\mathrm A)$ est libre.};}
% archive sheet number: 2
\scanpage{4}
\mspage{( 3 )}

\textit{Alors toute $\mathrm A$-algèbre intègre qui est localisée d’une $\mathrm A$-algèbre finie (i. e. $\mathrm A$-module de type fini) $\mathrm C$ est analytiquement réduite.}
\margin{\handwritten{$C\supset A$ !!!}}
% A curved pencil line leads from C to the upper note.

(Pour appliquer le lemme, on remarque que \tikz[baseline=(x.base)]{\node[draw,ellipse,inner sep=1pt] (x) {$h_1$};} et $h_2$ sont stables par localisation.)
\margin{\handwritten{faux}}
% The line from faux ends at the oval round h_1.

De l’existence de lemmes de normalisation [pour les séries convergentes, \textit{voir} (\textsuperscript{5}) par exemple], on déduit :

\textsc{Corollaire 1} (Chevalley). --- \textit{Toute localité intègre est analytiquement réduite.}

\textsc{Corollaire 2.} --- \textit{Soient $k$ un corps d’exposant caractéristique $p$ tel que $(k:k^p)<+\infty$, $\mathrm B$ une algèbre de séries convergentes sur $k$ et $\mathrm A$ un anneau intègre localisé d’un anneau quotient de $\mathrm B$. Alors $\mathrm A$ est analytiquement réduit.}

Dans le même ordre d’idées, notons la :

\textsc{Proposition 7.} --- \textit{Soit $\mathrm A$ un anneau local nœthérien analytiquement intègre et tel que $\mathrm K_{\widehat{\mathrm A}}$ soit séparable sur $\mathrm K_{\mathrm A}$. Alors la fermeture intégrale de $\mathrm A$ dans toute extension finie $\mathrm L$ de $\mathrm K_{\mathrm A}$ est un $\mathrm A$-module de type fini.}
% slanted
\margin{\handwritten{$\left.\rule{0pt}{2em}\right]$ connu (Nagata)}}

La proposition 7 s’applique aux localités, aux locali\uline{tés analytiques} et aux algèbres de Hopf régulières sur un corps parfait.
\margin{\handwritten{?}}

D’après le lemme 2 (p. 314, chap. 8, \S~13) de (\textsuperscript{11}), on voit alors par exemple qu’une localité analytique vérifiant les conditions du corollaire 2 et normale est analytiquement normale.
% slanted
\margin{\handwritten{$/$ connu}}

(\textsuperscript{*}) Séance du 4 novembre 1963.

(\textsuperscript{1}) \textsc{Berger} et \textsc{Kunz}, \textit{Math., Z.}, \textbf{77}, 1961, p. 314-338.

(\textsuperscript{2}) \textsc{Godement}, \textit{Séminaire Cartan-Chevalley}, 1955-1956, exp. 17, 18, 19.

(\textsuperscript{3}) \textsc{Kunz}, \textit{Math. Z.}, \textbf{76}, 1961, p. 56-74.

(\textsuperscript{4}) \textsc{Lang}, \textit{Ann. Math.}, \textbf{55}, 1952, p. 373-390.

(\textsuperscript{5}) \textsc{Mount}, \textit{Proc. Amer. Math. Soc.} juin 1963.

(\textsuperscript{6}) \textsc{Nagata}, \textit{Locals rings}, Inters., Tracts n\textsuperscript{o} 13, 1962.

(\textsuperscript{7}) \textsc{Nagata}, \textit{On the closedness of singular loci}, Publ. I. H. E. S., n\textsuperscript{o} 2, 1959.

(\textsuperscript{8}) \textsc{Nakai}, \textit{J. Math. Soc. Japan}, \textbf{13}, 1961, p. 63-84.

(\textsuperscript{9}) \textsc{Nakai} et \textsc{Suzuki}, \textit{J. Sc. Hiroshima Univ.}, série A, \textbf{23}, 1960, p. 459-476.

(\textsuperscript{10}) \textsc{Suzuki}, \textit{J. Math. Kyoto Univ.}, \textbf{2}, 1963, p. 157-182.

(\textsuperscript{11}) \textsc{Zariski} et \textsc{Samuel}, \textit{Commutative Algebra}, Van Nostrand.

\hfill(7, \textit{rue Chernoviz, Paris,} 16\textsuperscript{e}).
% archive sheet number: 3
\scanpage{5}
% Printer's imprint below a horizontal rule
\noindent\rule{\linewidth}{0.4pt}

\textbf{165.056. --- Imp. \textsc{Gauthier-Villars} \& C\textsuperscript{ie}, 55, Quai des Grands-Augustins, Paris (6\textsuperscript{e}).}

\centerline{\textbf{Imprimé en France.}}
% archive sheet number: 4
\scanpage{6}
\textit{C. R. Acad. Sc.}, t. \textbf{256}, p. 842-844.\hfill Séance du 21 janvier 1963.\hfill Groupe 1\par\noindent\rule{\linewidth}{0.4pt}

\textsc{Algèbre.} --- \textit{Sur les modules de différentielles.} Note (\textsuperscript{*}) de M. \textbf{\textsc{Jean-Pierre Jouanolou}}, présentée par M. René Garnier.

Application des modules de différentielles aux algèbres quasi-séparables et aux anneaux locaux réguliers.

$\mathrm A$ et $\mathrm B$ désignent dans toute cette Note deux anneaux commutatifs et unitaires, $u:\mathrm A\to\mathrm B$ un monomorphisme unitaire et $\mathrm D_{\mathrm A}(\mathrm B)$ le $\mathrm B$-module des différentielles de l’algèbre ainsi définie. Le passage de cette situation « globale » à la situation « locale » se fait en imposant à $\mathrm A$, $\mathrm B$, $u$ d’être locaux; l’idéal maximal de $\mathrm A$ (resp. $\mathrm B$) sera en ce cas noté $\mathfrak m$ (resp. $\mathfrak n$) et son corps résiduel $\mathrm L_a$ (resp. $\mathrm L_b$); par ailleurs, le rang de $\mathfrak m$ est $\operatorname{rg}\mathfrak m=\dim_{\mathrm L_a}(\mathfrak m/\mathfrak m^2)$, et de même pour $\mathfrak n$. Rappelons qu’étant donnée une extension de corps : $\mathrm K\to\mathrm L$, la dimension du $\mathrm L$-espace vectoriel $\mathrm D_{\mathrm K}(\mathrm L)$ représente en caractéristique $0$ (resp. $p>0$) le degré de transcendance de $\mathrm L$ par rapport à $\mathrm K$ (resp. le nombre d’éléments d’une $p$-base de $\mathrm L$ par rapport à $\mathrm K$) : on la notera $\{\mathrm L:\mathrm K\}$ dans tous les cas.

Nous nous proposons d’étudier d’abord les cas où $\mathrm D_{\mathrm A}(\mathrm B)$ est nul, puis ceux où il est libre (en liaison avec la régularité des anneaux locaux), le tout grâce aux deux suites exactes établies dans le cadre local par Berger et Kunz (\textsuperscript{1}) :
\[
\begin{aligned}
0\to\frac{\mathfrak n}{\mathfrak n^2+\mathfrak m\mathrm B}&\to\mathrm L_b\otimes_{\mathrm B}\mathrm D_{\mathrm A}(\mathrm B)\to\mathrm D_{\mathrm L_a}(\mathrm L_b)\to0,\\
0\to\frac{\mathfrak n}{\mathfrak n^2+(\mathfrak n\cap\mathrm A[\mathrm B^p])\,\mathrm B}&\to\mathrm L_b\otimes_{\mathrm B}\mathrm D_{\mathrm A}(\mathrm B)\to\mathrm D_{\mathrm L_a}(\mathrm L_b)\to0.
\end{aligned}
\]

La première est vraie si $\mathrm L_b$ est séparable sur $\mathrm L_a$, la deuxième, sinon : alors $p>0$ désigne la caractéristique des corps résiduels.

1. \textsc{Structure des algèbres quasi-séparables.} --- Une $\mathrm A$-algèbre $\mathrm B$ est dite \uline{quasi-séparable} lorsque $\mathrm D_{\mathrm A}(\mathrm B)=0$. Il s’agit là d’une propriété locale, en ce sens que $\mathrm B$ est quasi-séparable si pour tout idéal maximal $\mathfrak n$ de $\mathrm B$, $\mathrm B_{\mathfrak n}$ est quasi-séparable sur $\mathrm A_{\mathfrak m}$ ($\mathfrak m=\mathfrak n\cap\mathrm A$), et seulement en ce cas.
% slanted
\margin{\handwritten{\uline{non ramifiée}}}

\textit{Localement :}

(i) Si $\mathrm L_b$ est séparable sur $\mathrm L_a$ (en particulier en caractéristique $0$) : lorsque $\mathrm B$ est quasi-séparable,
\[
\mathfrak n=\mathfrak n^2+\mathfrak m\mathrm B,\tag*{(R)}
\]
et $\mathrm L_b$ est algébrique (resp. relativement parfait (\textsuperscript{2})) sur $\mathrm L_a$ en caractéristique $\mathrm O$ (resp. $p>0$) la réciproque étant vraie si l’on suppose, en outre, que $\mathrm D_{\mathrm A}(\mathrm B)$ est de type fini. Notons que si $\mathrm B$ est nœthérien, ou plus généralement si $\mathfrak n$ est de type fini, la relation (R) s’écrit
\[
\mathfrak n=\mathfrak m\mathrm B,
\]
ce qui exprime que $\mathrm B$ n’est pas ramifié sur $\mathrm A$.
\margin{\handwritten{!?}}
% archive sheet number: 5
\scanpage{7}
\mspage{( 2 )}

(ii) En caractéristique $p>0$ (pour les corps résiduels) : si $\mathrm B$ est quasi-séparable, $\mathrm L_b$ est relativement parfait (\textsuperscript{2}) sur $\mathrm L_a$ et
\[
\mathfrak n=(\mathfrak n\cap\mathrm A[\mathrm B^p])\,\mathrm B+\mathfrak n^2,\tag*{(S)}
\]
la réciproque étant vraie si l’on suppose en outre $\unsure{\mathrm D^{\mathrm A}|\mathrm D_{\mathrm A}}(\mathrm B)$ de type fini. Si $\mathrm B$ est nœthérien, (S) se simplifie encore.

\textit{Globalement :}

(i) Soit $\mathrm A$ un anneau nœthérien quasi-séparable sur un corps $\mathrm K$ :

--- en caractéristique $0$, $\mathrm A$ est composé direct d’une famille finie d’extensions algébriques de $\mathrm K$;

--- en caractéristique $p>0$, si pour tout idéal maximal de $\mathrm A$ le corps résiduel est une extension séparable (resp. de type fini) de $\mathrm K$, $\mathrm A$ est composé direct d’une famille finie d’extensions séparables et quasi-séparables (resp. algébriques séparables) de $\mathrm K$.

(ii) Soit $\mathrm A$ une algèbre quasi-séparable sur un corps $\mathrm K$ :

--- en caractéristique $0$, tout idéal premier de $\mathrm A$ est maximal et tous les corps résiduels des localisés de $\mathrm A$ sont des extensions algébriques de $\mathrm K$;

--- en caractéristique $p>0$, si pour tout idéal premier $\mathfrak g$ de $\mathrm A$ le corps $\mathrm A\mathfrak g/\mathfrak g\mathrm A\mathfrak g$ est de type fini sur $\mathrm K$, alors tout idéal premier de $\mathrm A$ est maximal et tous les corps résiduels sont des extensions algébriques séparables de $\mathrm K$;

--- Sous l’une ou l’autre de ces hypothèses, le quotient de $\mathrm A$ par sa nilracine est un anneau absolument plat (\textsuperscript{3}).

2. \textsc{Différentielles et anneaux locaux réguliers.} --- On se place dans le cadre local, en supposant même $\mathrm A$ et $\mathrm B$ intègres, de corps des fractions $\mathrm K_a$ et $\mathrm K_b$, et les cardinaux $\{\mathrm L_b:\mathrm L_a\}$, $\{\mathrm K_b:\mathrm K_a\}$, $\operatorname{rg}\mathfrak n$ finis.

\textsc{Théorème.} --- \textit{Si $\mathrm D_{\mathrm A}(\mathrm B)$ est de type fini (en particulier si $\mathrm B$ est de type fini sur $\mathrm A$) il y a équivalence entre :}

(i) \textit{$\mathrm D_{\mathrm A}(\mathrm B)$ est libre;}

(ii) \textit{a. Dans le cas où $\mathrm L_b$ est séparable sur $\mathrm L_a$ :}
\[
\dim_{\mathrm L_b}\frac{\mathfrak n}{\mathfrak n^2+\mathfrak m\mathrm B}=\{\mathrm K_b:\mathrm K_a\}-\{\mathrm L_b:\mathrm L_a\};
\]

\textit{b. Sinon, $p>0$ désignant la caractéristique de $\mathrm L_b$ et $\mathrm L_a$ :}
\[
\dim_{\mathrm L_b}\frac{\mathfrak n}{\mathfrak n^2+(\mathfrak n\cap\mathrm A[\mathrm B^p])\,\mathrm B}=\{\mathrm K_b:\mathrm K_a\}-\{\mathrm L_b:\mathrm L_a\}.
\]

Notons que l’hypothèse « $\mathrm D_{\mathrm A}(\mathrm B)$ de type fini » entraîne que les cardinaux $\{\mathrm K_b:\mathrm K_a\}$ et $\{\mathrm L_b:\mathrm L_a\}$ sont finis. De plus, l’égalité (ii) \textit{a} s’écrit :

--- Si $\mathrm A$ est un corps : $\operatorname{rang}\mathfrak n=\{\mathrm K_b:\mathrm A\}-\{\mathrm L_b:\mathrm A\}$,

--- Si $\mathrm B$ est non ramifié sur $\mathrm A$ : $\{\mathrm K_b\mid\mathrm K_a\}=\{\mathrm L_b\mid\mathrm L_a\}$.
% archive sheet number: 6
\scanpage{8}
\mspage{( 3 )}

\textsc{Corollaire.} --- \textit{Si $\mathrm A$ est une localité sur le corps $\mathrm K$, telle que $\mathrm K_a$ et $\mathrm L_a$ soient séparables sur $\mathrm K$, il y a équivalence entre :}

(i) \textit{$\mathrm D_{\mathrm K}(\mathrm A)$ est libre;}

(ii) $\mathrm A$ est régulier.

\textsc{Proposition 1.} --- \textit{Si $\mathrm A$ domine un corps $\mathrm K$, tel que $\mathrm D_{\mathrm K}(\mathrm A)$ soit libre et de rang fini, si $\mathrm K_b$ (resp. $\mathrm L_a$) est séparable sur $\mathrm K_a$ (resp. $\mathrm L_b$) et si}
\[
\operatorname{rg}\mathfrak n-\operatorname{rg}\mathfrak m=\{\mathrm K_b:\mathrm K_a\}-\{\mathrm L_b:\mathrm L_a\},
\]
\textit{$\mathrm D_{\mathrm A}(\mathrm B)$ est un $\mathrm B$-module de dimension homologique finie $\leq1$.}

\textsc{Corollaire.} --- \textit{Si $\mathrm A$ est un anneau local régulier contenant un corps $\mathrm K$, tel que le couple $(\mathrm K,\mathrm A)$ jouisse des propriétés énoncées dans le corollaire ci-dessus, $\mathrm D_{\mathrm A}(\mathrm B)$ est de dimension homologique $\leq1$ sous les hypothèses de la proposition 1.}

\textsc{Proposition 2.} --- \textit{Soit $\mathrm A$ un anneau local nœthérien contenant un corps $\mathrm K$ tel que $\{\mathrm L_a:\mathrm K\}$ soit fini et que $\mathrm L_a$ soit séparable sur $\mathrm K$. Si $\mathrm D_{\mathrm K}(\mathrm A)$ est libre, de rang $\dim\mathrm A+\{\mathrm L_a:\mathrm K\}$, $\mathrm A$ est local régulier.}

Notons, pour conclure, que dans de nombreuses questions (\textit{cf.} géométrie formelle) il vaut mieux séparer $\mathrm D_{\mathrm A}(\mathrm B)$ en considérant
\[
\mathrm D'_{\mathrm A}(\mathrm B)=\frac{\mathrm D_{\mathrm A}(\mathrm B)}{\displaystyle\bigcap_{n\geq0}\mathfrak n^n\mathrm D_{\mathrm A}(\mathrm B)};
\]
la proposition 2 est encore vraie, lorsqu’on y substitue $\mathrm D'_{\mathrm K}(\mathrm A)$ à $\mathrm D_{\mathrm K}(\mathrm A)$.

(\textsuperscript{*}) Séance du 14 janvier 1963.

(\textsuperscript{1}) \textsc{Berger} et \textsc{Kunz}, \textit{Math. Z.}, \textbf{77}, 1961, p. 314-338.

(\textsuperscript{2}) \textsc{Bourbaki}, \textit{Algèbre}, chap. V, \S~8, ex. 1.

(\textsuperscript{3}) \textsc{Bourbaki}, \textit{Algèbre comm.}, chap. I, \S~2, ex. 17.

(\textsuperscript{4}) \textsc{Godement}, \textit{Séminaire Cartan-Chevalley}, 1955-1956, exposés 17 et 18.

(\textsuperscript{5}) \textsc{Nakai}, \textit{J. Math. Soc. Japan}, 1961.

(\textsuperscript{6}) \textsc{Nakai} et \textsc{Suzuki}, \textit{J. Sc. Hiroshima Univ.}, série A, \textbf{24}, n\textsuperscript{o} 3, décembre 1960.

\hfill(\textit{École Normale Supérieure,} 45, \textit{rue d’Ulm, Paris,} 5\textsuperscript{e}.)
% archive sheet number: 7
\scanpage{9}
% Printer's imprint below a horizontal rule
\noindent\rule{\linewidth}{0.4pt}

\textbf{163.228. --- Imp. \textsc{Gauthier-Villars} \& C\textsuperscript{ie}, 55, Quai des Grands-Augustins, Paris (6\textsuperscript{e}).}

\centerline{\textbf{Imprimé en France.}}
% archive sheet number: 8
\scanpage{10}
\mspage{\handwritten{24}}

\margin{\handwritten{$/^{\circ}$ [\quad] inutile, résulte du \deleted{\unsure{fait|fait dit}} que $IB$ (\unsure{ou.|ou|resp.} $IC$) top. nilpotent.}}
\uline{contenant $\mathrm{IB}$ resp. $\mathrm{TC}$}\handwritten{$] /^{\circ}$}, enfin on suppose $\mathrm{IB}$ et $\mathrm{IC}$ fermés (condition automatiquement vérifiée si $\mathrm B,\mathrm C$ sont noethériens). \deleted{\unsure{Soit u:|Soit u}}Soient $\mathrm B_{\unsure{0|o}}=\mathrm B/\mathrm{IB}$,\deleted{et} $\mathrm C_0=\mathrm C/\mathrm{IC}$, $u:\inserted{\mathrm B\handwritten{\to}\mathrm C}$ un homomorphisme continu de $\mathrm A$-algèbres, $u_i:\mathrm B_i\handwritten{\to}\mathrm C_i$ l’homomorphisme \deleted{induit} déduit par passage au quotient (j’ai oublié d’introduire un système fondamental d’idéaux ouverts $\mathrm I_i\handwritten{\subset}\mathrm I_0\inserted{=\mathrm I}$ dans $\mathrm A$, et les $\mathrm B_i$ et $\mathrm C_i$ en conséquence\inserted{enfin ce sont tous les $\mathrm I_i\mathrm B$ et $\mathrm I_i\mathrm C$ qui doivent être fermés,}). Conditions équivalentes:

(i\handwritten{bis}) $u$ est \deleted{un isomorphisme} bijectif, et \tikz[baseline=(x.base)]{\node[draw,ellipse,inner sep=1pt] (x) {$u_0$ est un homéomorphisme.};}\handwritten{$\uparrow$}

\handwritten{$\curvearrowright$}(i\deleted{\unsure{bis}}) $u$ est un isomorphisme d’algèbres topologiques
% A curved arrow transposes this line to the position above (i bis).

(ii) Pour tout $i$, $u_i:\mathrm B_i\handwritten{\to}\mathrm C_i$ est bijectif, \tikz[baseline=(x.base)]{\node[draw,ellipse,inner sep=1pt] (x) {$u_0$ est un homéomorphisme.};}\handwritten{$\uparrow$}

(ii bis) pour tout $i$, muniss\deleted{\unsure{ons}}ant $\mathrm B_i$ et $\mathrm C_i$ des filtrations $\mathrm I$-adiques, l’homomorphisme $\operatorname{gr}_{\mathrm I}(\mathrm B_i)\handwritten{\to}\operatorname{gr}_{\mathrm I}(\mathrm C_i)$ induit par $u$\deleted{\unsure{e}} est bijectif, \tikz[baseline=(x.base)]{\node[draw,ellipse,inner sep=1pt] (x) {$u_0$ un homéomorphisme.};}\handwritten{$\uparrow$}
% slanted
\margin{\handwritten{$\bigcirc$ le \unsure{supposer} dès le début}}
% The ovals and upward pointers mark the repeated homeomorphism clauses.

Enfin, lorsque les $\deleted{\mathrm B}\handwritten{\mathrm C}_i$ sont plats sur les $\mathrm A_i$, \deleted{\unsure{les}}ces conditions équivalent également à

(iii) $u_0:\mathrm B_0\handwritten{\to}\mathrm C_0$ est un \deleted{isomorphisme\unsure{,}} \tikz[baseline=(x.base)]{\node[draw,ellipse,inner sep=1pt] (x) {homéomorphisme.};}\handwritten{$\uparrow$}

L’équivalence de (i) et (i bis) résulte trivialement du fait que les topologies de $\mathrm B,\mathrm C$ sont \deleted{adiques} définies par les puissances d’idéaux contenant $\mathrm{IB}$ resp. $\mathrm{IC}$. \deleted{Celle de} Evidemment (i\handwritten{bis})$\handwritten{\Rightarrow}$(ii), pour (ii)$\handwritten{\Rightarrow}$(i\handwritten{bis}) on note que les applications canoniques $\mathrm B\handwritten{\to}\underset{\handwritten{\longleftarrow}}{\lim}\mathrm B_i$ et $\mathrm C\handwritten{\to}\underset{\handwritten{\longleftarrow}}{\lim}\mathrm C_i$ sont bijectives grâce au fait que les $\mathrm I_i\mathrm B$ et $\mathrm I_i\mathrm C$ sont fermés (cf rectification à $0_{7.2.4.}$). L’équivalence de (ii) et (ii bis), \deleted{qu} mise pour mémoire, provient du fait que ($\mathrm I=\mathrm I_0$ étant un idéal de définition, donc nilpotent dans le\deleted{\unsure{s}}s $\mathrm A_i$) les filtrations mises sur $\mathrm B,\mathrm C$ sont finies. Enfin, reste à prouver que lorsque \deleted{\unsure{B est p}} $\deleted{\mathrm B}\handwritten{\mathrm C}_i$ est plat sur $\mathrm A_i$, alors $u_0$ bijectif implique $u_i$ bijectif ou si on veut $\operatorname{gr}(u)$ bijectif, ce qui est immédiat et sauf erreur figure déjà en lemme quelque part au Chap III (j’espère, avec la généralité qui convient, et bien entendu modulé ...). \handwritten{Cf \unsure{diagramme}}

\diagram{\handwritten{\begin{tikzcd}[ampersand replacement=\&]
\operatorname{gr} B \arrow[r,"\mathrm{surj}"] \& \operatorname{gr} C \\
B_0\otimes\operatorname{gr} A_i \arrow[u,"\mathrm{surj}"] \arrow[r,"\mathrm{bij}"'] \& C_0\otimes\operatorname{gr} A_{\unsure{i}} \arrow[u,"\mathrm{bij}"']
\end{tikzcd}}}
% archive sheet number: 9
\scanpage{11}
\mspage{\handwritten{25}}

Enfin, il nous sera commode d'avoir un critère permettant de déduire d\deleted{\unsure{u|a}}e la propriété $\mathrm P$ une propriété de platitude. Nous utilisero\deleted{j}ns :

\uline{Lemme} \deleted{\unsure{3|5}}\handwritten{4}.3. Soient $\mathrm A$ un anneau, $(\mathrm P_i)_{i\handwritten{\in}\mathrm I}$ une famille de $\mathrm A$-modules, \inserted{$\mathrm P$ leur prod\handwritten{uit}} $\mathrm M$ un $\mathrm A$-module de présentation finie. Alors l'homomorphisme canonique
\[
\mathrm M\otimes\mathrm P\handwritten{\to\prod}\mathrm M\otimes\mathrm P_i
\]
est un isomorphisme.

En effet, écrivant $\mathrm M$ comme \deleted{\unsure{un}}comme un conoyau d'un homomorphisme de modules libres \inserted{de type fini} $\mathrm L'\handwritten{\to}\mathrm L$, et notant que le foncteur produit est exact donc exact à gauche, on est ramené au cas où $\mathrm M=\mathrm A^n$, $n$ entier $\handwritten{>}0$, où c'est trivial.

\uline{Corollaire} \deleted{3}\handwritten{4}.4. Supposons $\mathrm A$ tel que tout module de type fini sur $\mathrm A$ soit de présentationfinie, par exemple $\mathrm A$ noethérien. Alors tout produit de $\mathrm A$-modules plats est un $\mathrm A$-module plat.

En effet, il suffit de montrer que $\mathrm P\otimes$ transforme un monomorphisme $\mathrm M\handwritten{\to}\mathrm N$ en un monomorphisme, quand $\mathrm M$ et $\mathrm N$ sont de type fini; le lemme 3 donne alors immédiatement le résultat voulu.

\uline{Corollaire} \deleted{3}\handwritten{4}.5. Supposons que $\mathrm A$ satisfasse \deleted{les conditions}\handwritten{la condition} de \deleted{3}\handwritten{4}.4.\inserted{par exemple $\mathrm A$ noethérien} Soit $\mathrm P$ un $\mathrm A$-module \deleted{topologique} isomorphe à \deleted{une}la limite projective \deleted{\unsure{``stricte'' dénomb}} \deleted{de $\mathrm A$-modules projectifs de type fini (``stricte'' = par un système projectif strict),} \deleted{``dénombrable'' par un système proj\unsure{e|x}} \deleted{Alors $\mathrm P$ est un $\mathrm A$} d'un système projectif strict,\deleted{de $\mathrm A$} admettant un système cofinal dénombrable, de $\mathrm A$-modules projectifs. Alors $\mathrm P$ est un $\mathrm A$-module plat.

En effet, on voit faci\deleted{\unsure{u}}lement que sous les conditions dites, $\mathrm P$ est isomorphe à un produit de $\mathrm A$-modules projectifs, donc plats, d'où le résultat.
% archive sheet number: 10
\scanpage{12}
\mspage{\handwritten{26}}

\uline{Théorème} \deleted{\unsure{3}}\handwritten{4}.6. Soient $\mathrm A$ un anneau préadmissible muni d'un idéal de définition $\mathrm I=\mathrm I_{\unsure{0|o}}$ et d'un système fondamental $(\mathrm I_i)_i$ d'idéaux d e définition $\mathrm I_0$, $\mathrm B$ une $\mathrm A$-algèbre topologique. On suppose\deleted{\unsure{r}} les $\mathrm A_i=\mathrm A/\mathrm I_i\mathrm A$ noethériens, et de plus les conditions suivantes sur $\mathrm B$:
\deleted{\unsure{a) Les $\mathrm A_i=\mathrm A/\mathrm I_i$ noethériens|a) les $\mathrm A_i$ sont nœthériens}}

\deleted{\unsure{b}}\handwritten{a}) Les \deleted{\unsure{B|M}}\handwritten{m}odule\deleted{s} topologique\deleted{s} \deleted{$\unsure{\mathrm B_i=\mathrm B/\mathrm I_i\mathrm B}$} \deleted{\unsure{sur} $\unsure{\mathrm A_i}$} \inserted{$\mathrm B$ sur $\mathrm A$} satisf\deleted{on}\handwritten{ai}t la condition $\mathrm P$.

\inserted{\handwritten{b) $B_0$ est formellement simple sur $A_0$}}

\deleted{b}\handwritten{c}) $\mathrm B$ est adique \deleted{pour un idéal contenant $\mathrm{IB}$}

\handwritten{d}) Les idéaux $\mathrm I_i\mathrm B$ dans $\mathrm B$ sont fermés (\deleted{inutile}\inserted{\handwritten{\unsure{automatiquement} vérifié}} si $\mathrm B$ est noethérien).

Sous ces conditions, on a ce qui suit\deleted{\unsure{e}}:

(i) $\mathrm B$ est formellement simple sur $\mathrm A$.

(ii) Soit $\mathrm C$ une $\mathrm A$-algèbre topologique satisfaisant à \deleted{\unsure{b}}\handwritten{c}) et \handwritten{d}) (mais pas nécessairement à a) ), et soit $u_0:\mathrm B_0\handwritten{\overset{\sim}{\to}}\mathrm C_0$ un isomorphisme d'al\`{g}èbres topologiques. Alors il existe un $\mathrm A$-homomorphisme continu $u:\mathrm B\handwritten{\to}\mathrm C$ se réduisant suivant $u_0$ \handwritten{(}en vertu de (i) et de 1.7.). Ce dernier étant choisi, les conditions suivantes sont équivalentes: \deleted{\unsure{ce}}

\deleted{\unsure{a)}}1) $\mathrm B$ est formellement simple sur $\mathrm A$.

2) \deleted{$\mathrm B$ est} $\mathrm B$ satisfait la c\deleted{\unsure{o}}ondition a) ci-dessus, i.e. \deleted{est} satisf\deleted{\unsure{i|.}}ait la condition $\mathrm P$.

3) Les $\mathrm B_i$ sont plats sur $\mathrm A_i$.

4) L'homomorphisme $u:\mathrm B\handwritten{\to}\mathrm C$ est un isomorphisme\deleted{\unsure{,}} de $\mathrm A$-algèbres top.

\inserted{\handwritten{5) $\deleted{\unsure{B}}C$ est isom. à $B$ comme $A$-algèbre top.}}

Démonstration. L'assertion (i) résulte déjà de a)\inserted{\handwritten{ et b)}}, grâce à 3.1.

Prouvons donc (ii). En vertu de (i) d\deleted{\unsure{q}}éjà prouvé, on a 2)$\handwritten{\Rightarrow}$1), et \deleted{\unsure{e}}Evidemment 4)$\handwritten{\Rightarrow}$2), d'ailleurs 2)$\handwritten{\Rightarrow}$3) en vertu de 3.5. , grâce au fait que les $\mathrm A_i$ sont noethériens. Reste à prouver 1)$\handwritten{\Rightarrow}$4) et 3)$\handwritten{\Rightarrow}$4).

Supposons 1) vérifié, alors en vertu de 1.7. $u_0^{\handwritten{-}1}:\mathrm C_0\handwritten{\to}\mathrm B_0$ se prolonge en un homomorphisme d'algèbres top $v:\mathrm C\handwritten{\to}\mathrm B$ (grâce aux hypothèses \deleted{\unsure{b}}\handwritten{c}) et \handwritten{d}) ). Or on a déjà noté que les $\mathrm B_i$ sont plats sur les $\mathrm A_i$. Le lemme \deleted{3}\handwritten{4}.2. (i) et (iii) (où on échange les rôles de $\mathrm B,\mathrm C$)
% archive sheet number: 11
\scanpage{13}
\mspage{\handwritten{27}}

garantit alors que $v$ est un isomorphisme d'algébres topologiques (grâce à \deleted{\unsure{b}}\handwritten{c}) et \handwritten{d}) pour $\mathrm B$ et $\mathrm C$). Donc \deleted{les} on a 2), \inserted{donc on a 3). Or} \deleted{et} le même argument en sens inverse montre que \deleted{$u\unsure{{:}\mathrm C\to\mathrm B}$ est alors un isomorphisme d'alg top, i.e. qu'on a 4).} \inserted{3) implique 4). - Enfin 4)$\handwritten{\Rightarrow}$5) et 5)$\handwritten{\Rightarrow}$2).} On peut certainement améli\deleted{\unsure{r}}orer la rédaction de la démonstratio\deleted{\unsure{s}}n, évitant certaines redites.

Pour simplifier, nous nous restreindrons par la suite à des hypothèses noethériennes (permettant de nous déb\deleted{\unsure{aa|rr}}\unsure{ar|ra|rr|arr}asser en particulier de la désagréable hypothèse \deleted{\unsure{c}}\handwritten{d}) sur $\mathrm B$).

\uline{Définition} \deleted{\unsure{3|5}}\handwritten{4}.7. Soit $\mathrm A\handwritten{\to}\mathrm B$ un homomorphisme continu d'anneaux noethériens préadiques. On dit que $\mathrm B$ est une \uline{algèbre} (topologique) \uline{de Cohen} sur $\mathrm A$ si \inserted{\handwritten{a)}} $\mathrm B$ est séparé et complet, i.e. \inserted{\handwritten{est}} un anneau adique,\deleted{sur $\mathrm A$,} \inserted{\handwritten{b)}} s'il satisfait à la condition $\mathrm P$, et si enfin \inserted{\handwritten{c)}} $\mathrm B$ est formellement \deleted{\unsure{plat}}\inserted{simple} sur $\mathrm A$ \deleted{\unsure{On}}\inserted{\handwritten{, (ou encore, $I$ étant un \unsure{idéal} de \unsure{définition} de $A$, $B_{\unsure{0|o}}=B/IB$ est form\textsuperscript{t} simple sur $A_0=A/I$)}}

On notera qu'en vertu de 3.1. , si $\mathrm I$ est un idéal de définition de $\mathrm A$, la dernière condition équivaut à la condition que $\mathrm B/\mathrm{IB}$ est formellement simple sur $\mathrm A/\mathrm{IA}$. D'ailleurs, on peut interpréter \deleted{3}\handwritten{4}.6. de la façon suivante (et il est \deleted{peutêtre} plus raisonnable de réléguer \deleted{3}\handwritten{4}.6. en lemme, et de réserver le corollaire en théorème):

\uline{Coroll\deleted{\unsure{e}}aire} \deleted{3}\handwritten{4}.8. Soient $\mathrm A$ un anneau \deleted{préadique} noethérien, $\mathrm I$ un idéal de $\mathrm A$, munissons $\mathrm A$ de la topologie $\mathrm I$-préadique, posons $\mathrm A_0=\mathrm A/\mathrm{IA}$, et soit $\mathrm B_0$ une $\mathrm A_0$-algèbre topologique. \deleted{S'il}\inserted{Supposons qu'\handwritten{il}} existe une algèbre de Cohen $\mathrm B$ sur $\mathrm A$ telle que $\mathrm B/\mathrm{IB}$ soit isomorphe à $\mathrm B_0$ comme $\mathrm A_0$-algèbre topologique,\inserted{(ce qui implique $\mathrm B_0$ \deleted{\unsure{formt simple}}\handwritten{anneau de Cohen} sur $\mathrm A_0$).} alors $\mathrm B$ est unique à \inserted{un} isomorphisme près \deleted{(induisant l'isomorphisme donné sur} (compatible avec les isomorphismes donnés des algèbres \deleted{\unsure{réduite|réduites}} réduites mod $\mathrm I$ avec $\mathrm B_0$). De plus, si $\mathrm B$ est une $\mathrm A$-algèbre topologique adique telle que $\mathrm B/\mathrm{IB}$ soit \deleted{\unsure{A}}isomorphisme à $\mathrm B_0$ comme $\mathrm A_0$-algèbre topologique, alors les conditions suivantes sur $\mathrm B$ sont équivalentes: \handwritten{\unsure{0|a})} $\mathrm B$ est une $\mathrm A$-algèbre de Cohen
% archive sheet number: 12
\scanpage{14}
\mspage{\handwritten{28}}

1) $\mathrm B$ est formellement simple sur $\mathrm A$

\deleted{\unsure{3}}2) $\mathrm B$ satisfait la condition $\mathrm P$\quad 3) $\mathrm B$ est plat sur $\mathrm A$\quad 3 bis) Pour tout entier $n\handwritten{\geq}0$, $\mathrm B_n=\mathrm B/\mathrm I^{n+1}\mathrm B$ est plat sur $\mathrm A_n=\mathrm A/\mathrm I^{n+1}$. \deleted{4)}

Il reste seulement à montrer l'équivalence de 3) et 3 bis), qui est contenue en effet (moyennant les hypothèses noethériennes faites) dans la platitude supérieure.

\deleted{Le résultat précédent} Il conviendrait de noter aussi\inserted{ dès après 3.7.} que ``algèbre de Cohen'' est une notion stable par extension de la base via produits tensoriels complétés: c'est une trivialité.

Le résultat \deleted{\unsure{3}}\handwritten{4}.8. montre l'intérêt qu'il y a, $\mathrm A,\mathrm I$ et $\mathrm B_{\unsure{0|o}}$ étant donnés, de disposer d'un \uline{théorème d'existence} pour une algèbre de Cohen $\mathrm B$ sur $\mathrm A$ se réduisant suivant $\mathrm B_0$. On notera que la construction habituelle de proche en proche nous ramène \deleted{au cas où $\mathrm A$ est discr}(remplaçant $\mathrm A$ par des $\mathrm A_i$ et $\mathrm I$ par des puissances) au cas où $\mathrm I$ est de carré nul,(donc où $\mathrm A$ est discret). On obtient alors un problème de nature essentiellement infinitésimale, dont nous ignorons s'il existe toujours une solution, faute de savoir caractériser la structure des algèbres de Cohen $\mathrm B_0$ sur l'anneau discret $\mathrm A_0$. Cependant, nous obtenons \inserted{\handwritten{(\deleted{\unsure{3}}\unsure{4|3}.10.)}} un théorème d'existence dans un cas particulier intéressant, en utilisant les résultats du N\textsuperscript{\uline{o}}2\deleted{\unsure{,}}. Notons \inserted{d'abord}

\deleted{\uline{Théorème} 3.9. Soit $\mathrm A$ \deleted{$\mathrm B$ un homomorphisme} un anneau local noethérien, \deleted{et de corps} résiduel $k$, et soit $\mathrm B_0$ une $k$-algèbre, \deleted{qui soit \unsure{u|n} satisfasse} \inserted{\deleted{qui} soit} \deleted{aux conditions suivantes: $\mathrm B_0$ est} un anneau local noethérien, \deleted{et complet,} \inserted{de} \deleted{son corps} résiduel $\mathrm K$ \deleted{est de degré de multiplicité radicielle finie sur $k$ (ce qui est le cas} une extension séparable de $k$, $\mathrm B_0$ est (pour sa topolo-}

\uline{Corollaire} \deleted{3}\handwritten{4}.9. Soit $\mathrm A\handwritten{\to}\mathrm B$ un homomorphisme local d'anneaux locaux noethériens, \deleted{tel que l'extension résiduelle $k(\mathrm B)/k(\mathrm A)$ soit à multiplicité}
% archive sheet number: 13
\scanpage{15}
\mspage{\handwritten{29}}

\deleted{radicielle finie (par exemple, soit séparable)}. \inserted{\handwritten{Conditions \unsure{équivalentes}: 1)}} \deleted{Pour que} $\mathrm B$ \deleted{soit}\inserted{\handwritten{est}} une algèbre de Cohen sur $\mathrm A$ (sous-entendu: pour les topologies habituelles, - il faudrait le préciser dans \deleted{\unsure{3}}\handwritten{4}.7.) , \deleted{il f et suffit que}\inserted{\handwritten{2)}} $\mathrm B$ \deleted{\unsure{soit|est}}\handwritten{est} plat sur $\mathrm A$, complet, et \deleted{que} $\mathrm B_{\unsure{0|o}}=\mathrm B/\underline{\mathrm m}\mathrm B$ \deleted{soit}\inserted{\handwritten{est}} une algèbre de Cohen sur $k=\mathrm A/\underline{\mathrm m}$, i.e. formellement simple sur $k$, \inserted{\handwritten{cf 2.\unsure{6|5}.}} \deleted{\handwritten{3) $B$ est \inserted{\unsure{complet et}} formellement simple sur $A$ \unsure{à fibre|et tel que|ou plat et} \unsure{séparable}}}

(On aura noté \deleted{pa} après \deleted{3}\handwritten{4}.7. que pour une algèbre topologique sur un corps $k$, Cohen = adique et formellement simple ). \deleted{La nécessité}\inserted{\handwritten{1)$\Rightarrow$2) \deleted{\unsure{et 1)$\Rightarrow$3)}}}} étant déjà acquis\deleted{e} (sans hypothèse de multiplicité radicielle finie), prouvons \deleted{la suffisance}\inserted{\handwritten{2)$\Rightarrow$1)}}. En vertu de la définition, on est ramené simplement à prouver que $\mathrm B$ satisfait la condition $\mathrm P$. Cela nous ramène au cas où $\mathrm A$ est artinien. En vertu du N\textsuperscript{\uline{o}}2, on sait que $\mathrm B_0$ est régulier, considérons un système régulier de paramètres de $\mathrm B_0$ (un système de paramètres formant une $\mathrm B_0$-suite suffirait), relevons le en un système de paramètres de $\mathrm B$, soit $x=(x_1,\ldots,x_r)$. Alors les idéaux engendrés par les $x^n=(x_1^n,\ldotp\ldotp,x_r^n)$, pour $n\handwritten{\geq}0$, forment un système fondamental de voisinages \inserted{\handwritten{$\unsure{I_n|J_n}$}} de 0 \deleted{\handwritten{\unsure{de}}} dans $\mathrm B$, d'autre part on constate aussitôt que les $\mathrm B/\mathrm J_n$ sont des modules plats sur $\mathrm A$ (il y a un \deleted{lemme} corollaire dans l\deleted{\unsure{q}}a platitude supérieure, dans le cas où on divise par un seul élement \deleted{non divi} régulier, cas auquel on est ramené aussitôt). Comme $\mathrm A$ est artinien, ils sont donc projectifs. \deleted{\unsure{c.q.f.d.}} \deleted{\handwritten{\unsure{Enfin, on aura} 3)$\Rightarrow$2) \unsure{si $k(\mathrm B)/k$ séparable} \unsure{grâce à} 3.\unsure{8}. \unsure{(passé de 3) à 1))} en utilisant le \uline{th. d'existence} \uline{ci-dessous} 3.\unsure{9|11}.}}

\deleted{NB\quad On n'a utilisé l'hypothèse sur la multiplicité radicielle que pour \deleted{pouvoir} savoir que $\mathrm B_0$ est Cohen-Macaulay. L'hypothèse faite est très probablement inutile. Serre vient de me dire que Mac-Lane a prouvé la réciproque de 2.3. sans hypothèse de mu\deleted{\unsure{ll}}ltiplicité rad finie, dans ``Note on a relative structure of p-adic fi\deleted{\unsure{el}}elds'' Ann Math t.41 1940, page 751 -753. Cela vaudrait peut-être le coup, pour \deleted{liquider d} le coup d'oeuil, que tu regardes cela et donnes le résultat complet dans 2.3.}\margin{\handwritten{$\Vert$}} On
% A double vertical stroke stands beside the citation and page number. The ruled cancellation steps around the last On.
% archive sheet number: 14
\scanpage{16}
\mspage{\handwritten{30}}

\deleted{\deleted{pourra alors}\inserted{\handwritten{peut donc}} \inserted{dans 3.9.} resper l'hypothèse \deleted{dans 3.9.} avec : ``ou \deleted{\unsure{$\mathrm B$|$\mathrm B_0$} \unsure{est un anneau un corp|est un anneau un corps}} $\mathrm r(\mathrm B)=\mathrm r(\mathrm A)\mathrm B$'' \deleted{(car grâce au résultat de et la}}

\uline{Corollaire} \deleted{\unsure{3|5}}\handwritten{4}.10. Soi\deleted{t}ent $\mathrm A$ un anneau local noethérien, de corps résiduel $k$, $\mathrm B_{\unsure{0|o}}$ une $k$-algèbre qui soit un anneau local noethérien, et soit une $k$-algèbre de Cohen i.e. une $k$-algèbre formel\deleted{\unsure{h|l}}lement simple \deleted{pour sa} et complète pour sa topologie ordinaire. Supposons de plus l'extension résiduelle $\mathrm K/k$ séparable (ce qui n'est donc pas une restriction si $\mathrm B_0=\mathrm K$\deleted{\unsure{=k|,}}, mais en est évidemment une en général). Alors il existe une $\mathrm A$-algèbre topologique de Cohen \inserted{$\mathrm B$} qui se réduit suivant $\mathrm B_0$, $\mathrm B$ est unique à isomorphisme près en vertu de \deleted{3}\handwritten{4}.8. , et c'est un anneau noethérien local complet comme il se doit.

En vertu du théorème de structure \deleted{\unsure{2}}\handwritten{3}.4. on a $\mathrm B_0\handwritten{\simeq}\mathrm K\handwritten{[\![}\mathrm T_1,\ldotp\ldotp,\mathrm T_n\handwritten{]\!]}$, et on est ramené, grâce au critère \deleted{3}\handwritten{4}.9. 2) et la transitivité de la platitude, aux deux c\deleted{\unsure{q}}as types: 1\textsuperscript{\uline{o}}) $\mathrm B_0=\mathrm K$ 2\textsuperscript{\uline{o}}) $\mathrm B_0=k\handwritten{[\![}\mathrm T_1,\ldotp\ldotp,\mathrm T_n\handwritten{]\!]}$. Dans le premier cas, on utilise le résultat général de $0_{\mathrm{III}}$\inserted{\handwritten{10.3.1.}}, et dans le deuxième cas, il suffit de prendre $\mathrm B=\handwritten{\widehat{\mathrm A}[\![}\mathrm T_1,\ldots,\mathrm T_n\handwritten{]\!]}$.

\deleted{Théorème 3.1} \uline{Corollaire} \deleted{3}\handwritten{4}.11. Sous les conditions de 3.9. , \inserted{\handwritten{si $\unsure{K|k}(B)/k(A)$ de mult. rad. finie,}} les conditions 1) et 2) équivalent aussi à: 3) $\mathrm B$ est complet et formellement simple sur $\mathrm A$

En effet, on sait déjà que \deleted{c'est nécessaire,}\inserted{1)$\handwritten{\Rightarrow}$3)}, reste à prouver \deleted{la} 3)$\handwritten{\Rightarrow}$1) \deleted{suffisance} ou encore 3)$\handwritten{\Rightarrow}$2), i.e. on est ramené à prouver simplement que $\mathrm B$ est plat sur $\mathrm A$. Or\deleted{\unsure{n}} faisant une extension de la base $\mathrm A\handwritten{\to}\mathrm A'$, $\mathrm A'$ local libre \inserted{de type fini} sur $\mathrm A$ à extension radicielle convenable (cf $0_{\mathrm{III}}$), on trouve un $\mathrm B'$ \inserted{local} sur $\mathrm A'$ qui sera encore formellement simple sur $\mathrm A'$, et de plus à extension résiduelle \deleted{triviale} séparable\handwritten{$/^{\circ}$}. On dispose alors du théorème d'existence \deleted{3}\handwritten{4}.10. qui nous garantit, grâce à \deleted{3}\handwritten{4}.8. (équiv de 0) et 1) ) que $\mathrm B'$ est Cohen sur $\mathrm A'$donc plat sur $\mathrm A'$. Donc $\mathrm B$ est plat sur $\mathrm A$.

\margin{\handwritten{($/^{\circ}$ réduction inutile dans le cas ``\unsure{vespé|visé|vrai}'' du NB)}}
% archive sheet number: 15
